% RQ2 Introduction: Governance Capture Mitigation

\subsection{Motivation}

The promise of DAOs rests on distributed decision-making---governance by the many rather than the few. Yet in practice, most DAOs exhibit extreme power concentration. Studies of MakerDAO found that 20 addresses controlled over 50\% of voting power \citep{fritsch2022analyzing}. Similar patterns appear across major protocols: a handful of whales and delegates routinely determine proposal outcomes.

This concentration creates multiple risks:
\begin{itemize}
    \item \textbf{Plutocracy}: Wealthy actors dominate governance regardless of community preference
    \item \textbf{Capture}: Coordinated groups can manipulate outcomes for private benefit
    \item \textbf{Legitimacy erosion}: Small token holders disengage, viewing governance as futile
    \item \textbf{Centralization}: De facto control contradicts the decentralization thesis
\end{itemize}

Token-weighted voting---one token, one vote---directly maps economic stake to political power. While this aligns incentives (those with most to lose have most voice), it also enables capture by anyone willing to acquire sufficient tokens.

\subsection{Research Question}

This paper investigates:

\begin{quote}
\textbf{RQ2:} How effective are different mitigation strategies (vote caps, quadratic thresholds, velocity penalties) against governance capture?
\end{quote}

Specifically, we examine:
\begin{enumerate}
    \item How much do vote caps reduce whale influence, and at what cost to governance throughput?
    \item Does quadratic voting meaningfully redistribute power, or do whales simply split holdings?
    \item Can velocity penalties prevent rapid accumulation attacks without hindering legitimate participation?
    \item What combinations of mechanisms provide robust capture resistance?
\end{enumerate}

\subsection{Mitigation Strategies}

We evaluate three classes of capture mitigation:

\subsubsection{Vote Caps}

Hard limits on maximum voting power per address:
\begin{equation}
w_{\text{capped}}(a) = \min(w(a), c)
\end{equation}

This directly bounds individual influence but creates incentives for address splitting (sybil attacks).

\subsubsection{Quadratic Voting}

Diminishing returns on token power:
\begin{equation}
w_{\text{quad}}(a) = \sqrt{T(a)}
\end{equation}

This smoothly reduces whale influence without hard cutoffs, but the square root relationship may still permit substantial concentration.

\subsubsection{Velocity Penalties}

Reduced voting power for recently-acquired tokens:
\begin{equation}
w_{\text{velocity}}(a) = T_{\text{old}}(a) + \alpha \cdot T_{\text{new}}(a)
\end{equation}

where $\alpha < 1$ penalizes recent acquisitions. This addresses ``governance attacks'' where actors buy tokens immediately before votes.

\subsection{Contributions}

This paper contributes:

\begin{enumerate}
    \item \textbf{Systematic comparison} of three capture mitigation mechanisms under controlled conditions
    \item \textbf{Quantified tradeoffs} between capture resistance and governance efficiency
    \item \textbf{Optimal parameter ranges} for each mechanism derived from simulation
    \item \textbf{Hybrid recommendations} combining multiple strategies for robust defense
\end{enumerate}

\subsection{Paper Organization}

Section~\ref{sec:background} reviews governance capture in DAOs and existing defenses. Section~\ref{sec:theory} formalizes capture metrics and mitigation mechanisms. Section~\ref{sec:architecture} describes simulation implementation. Section~\ref{sec:methodology} details experimental design. Section~\ref{sec:results} presents findings. Section~\ref{sec:discussion} interprets results and derives practical guidance. Section~\ref{sec:conclusion} concludes.
